% !PRINTED 623
\weberpage{623}
% No page number is visible within the captured band crop (the top
% band is blank above the Abschnitt title, as this chapter opener carries no
% running head, and the foot of the page below the footnote is blank too). 623 is inferred from direct pagination continuity: this page's footnote
% 1) is continued verbatim at the head of the footnote apparatus on the next page,
% whose own running head clearly prints "624".

\begin{center}
\gesperrt{Vierundzwanzigster Abschnitt.}

\medskip
{\Large Algebraische Funktionen einer Variablen.}
\end{center}

\rule{5em}{0.4pt}

\begin{center}
§ 170. \gesperrt{Einleitendes.}
\end{center}

Die Theorie der algebraischen Funktionen einer Veränderlichen ist der
Ausgangspunkt der allgemeinen Untersuchungen von \gesperrt{Abel} über
die neuen Transzendenten, die seitdem den Namen \glqq \gesperrt{Abelsche
Funktionen}\grqq{} erhalten haben, und die höchste Verallgemeinerung
der elliptischen Funktionen sind. Die Hauptprobleme dieser Theorie
sind durch die Arbeiten von \gesperrt{Riemann}, \gesperrt{Weierstrass},
\gesperrt{Clebsch}, \gesperrt{Brill} und \gesperrt{Noether} zu einem
gewissen Abschluß gebracht. Insbesondere hat \gesperrt{Riemann} in der
Vorstellung der mehrblätterigen (\gesperrt{Riemann}schen) Flächen ein durch seine
Anschaulichkeit außerordentlich wirksames Hilfsmittel für die
Untersuchung dieser Funktionen geschaffen.

Alle diese Untersuchungen aber, die sich einerseits auf die
Funktionentheorie, andererseits, wie bei \gesperrt{Clebsch},
\gesperrt{Brill}, \gesperrt{Noether}, auf die Methoden der rationalen
Algebra (Theorie der Formen und Invarianten oder der algebraischen
Kurven) stützen, müssen immer gewisse Einschränkungen machen, sie
müssen gewisse \glqq Ausnahmefälle\grqq{} ausschließen und sich mit
der Behandlung der sogenannten allgemeinen Fälle begnügen. Nicht
unterworfen ist dieser Beschränkung die nach Analogie der
Zahlentheorie von \gesperrt{Dedekind} und mir ausgearbeitete Theorie
der algebraischen Funktionen, von der hier eine Übersicht gegeben
werden soll\textsuperscript{1)}.

\smallskip
\noindent\rule{5em}{0.4pt}

\noindent\textsuperscript{1)}~Zur Literatur über algebraische und
\gesperrt{Abel}sche Funktionen sei hier erwähnt: \gesperrt{Abel}, Mém. sur une
propriété générale d'une classe très-étendue de fonctions
transcendentes (1826 der Pariser Akademie vorgelegt, Werke, neue
% [sic] printed "transcendentes", Abel's French title (and French spelling) has "transcendantes"
Ausgabe von \gesperrt{Sylow} und \gesperrt{Lie}, Bd.\,I, S.\,145).
\gesperrt{Riemann}, Theorie der \gesperrt{Abel}schen Funktionen, Crelles Journ.,
Bd.\,54, 1857, Werke 2. Aufl., Nr.\,88. \gesperrt{Weierstrass},
Vorlesungen 1875/76, Werke, Bd.\,IV. \gesperrt{Clebsch}, Über die
%CONT>>
